\Askhsh{34408}{2}
\begin{ylh}{1}
    \item 4.6. Άθροισμα γωνιών τριγώνου
    \item 5.3. Ορθογώνιο
    \item 5.9. Μια ιδιότητα του ορθογώνιου τριγώνου.
\end{ylh}
% ===================================================================
%  Αρχείο: 34408.tex (Fragment)
%  Αυτόματη μετατροπή μέσω doc2latex.py
% ===================================================================

ΘΕΜΑ 2

Δίνεται τετράπλευρο ΑΒΓΔ με $\widehat{Α}$ = $\widehat{\Delta}$ = 90\textsuperscript{0}, ΑΒ \textgreater{} ΓΔ, ΒΓ = 4 ΔΓ και $\widehat{Β}$ = 60\textsuperscript{0} και ΓΗ Ʇ ΑΒ.

Να αποδείξετε ότι:

\begin{alist}
\item ΗΒ = 2ΔΓ, \monades{12}

\item το τετράπλευρο ΑΗΓΔ είναι ορθογώνιο με ΑΗ = $\frac{1}{2}$ ΗΒ. \monades{13}


\begin{center}
\begin{tikzpicture}[scale=1]
\tkzDefPoint(0,4){D}
\tkzDefPoint(0,0){A}
\tkzDefPoint(2,4){C} % Just to have a length CD = 2
\tkzDefPoint(6,0){B} % Just to have AB > CD. Let's make it consistent with problem.
% BC = 4 CD. If CD=x, BC=4x. Angle B = 60.
% CD is top base, AB is bottom. D, A are right angles.
% H is projection of C onto AB. AH = DC. HB = AB - DC.
% In right triangle CHB, angle B= 60. CH = HB * tan 60. BC = HB / cos 60 = 2 HB.
% Problem says BC = 4 CD. So 2 HB = 4 CD => HB = 2 CD.
% So we need HB = 2 CD.
\tkzDefPoint(0,0){A}
\tkzDefPoint(0,3.46){D} % height
\tkzDefPoint(1.5,3.46){C} % CD = 1.5
\tkzDefPoint(1.5,0){H}
\tkzDefPoint(4.5,0){B} % HB = 3 = 2 * CD.

\draw[l3] (A)--(B)--(C)--(D)--cycle;
\draw[l3] (C)--(H);

\tkzDrawPoints(A,B,C,D,H)
\tkzLabelPoint[below left](A){$A$}
\tkzLabelPoint[below right](B){$B$}
\tkzLabelPoint[above right](C){$\Gamma$}
\tkzLabelPoint[above left](D){$\Delta$}
\tkzLabelPoint[below](H){$H$}

\tkzMarkRightAngle(D,A,B)
\tkzMarkRightAngle(A,D,C)
\tkzMarkRightAngle(C,H,B)
\tkzMarkAngle[size=0.6,mark=none](C,B,H)
\tkzLabelAngle[pos=0.9](C,B,H){$60^\circ$}
\tkzMarkSegments[mark=s|](D,C A,H)
\tkzMarkSegments[mark=s||](H,B) % HB = 2 CD. Maybe better mark?
\end{tikzpicture}
\end{center}
\end{alist}
